\documentclass[11pt]{article}
\usepackage{amsmath,amssymb,amsthm}
\setlength{\textwidth}{6.3in}
\setlength{\oddsidemargin}{0.1in}
\setlength{\evensidemargin}{0.1in}
\setlength{\textheight}{8.6in}
\setlength{\topmargin}{-0.3in}
\setlength{\headheight}{14pt}
\setlength{\headsep}{22pt}
\setlength{\parindent}{1.2em}
\setlength{\parskip}{3pt}
\newcommand{\R}{\mathbb R}
\newcommand{\N}{\mathbb N}
\newcommand{\Nzero}{\mathbb N_0}
\newcommand{\gra}{\operatorname{gra}}
\newcommand{\dom}{\operatorname{dom}}
\newcommand{\sgn}{\operatorname{sgn}}
\newcommand{\pair}[2]{\langle #1,#2\rangle}
\newtheorem{theorem}{Theorem}
\newtheorem{proposition}{Proposition}
\newtheorem{lemma}{Lemma}
\newtheorem{example}{Example}
\newtheorem{corollary}{Corollary}
\theoremstyle{definition}
\newtheorem{definition}{Definition}
\newtheorem{assumption}{Assumption}
\theoremstyle{remark}
\newtheorem{remark}{Remark}
\newcommand{\dist}{\operatorname{dist}}
\pagestyle{plain}
\begin{document}

\begin{center}
{\Large\bfseries Nonmaximal sums of maximally monotone operators under Rockafellar’s constraint qualification}\\[10pt]
\begin{tabular}{c}
Weifeng Yang\footnotemark
% Yunnan Earthquake Agency, Kunming, China\\
\end{tabular}\\[10pt]
% September 7, 2026\\[4pt]
 % Main manuscript
\end{center}
\footnotetext{\texttt{Email:ywf841673182@gmail.com}}


\begin{abstract}


We construct counterexamples to Rockafellar's sum conjecture in which two maximally monotone operators satisfy the interior-domain condition but their sum is not maximally monotone. We give one counterexample on $c_0$ and another on $\ell^1$ with its usual norm.
We establish a general construction theorem that computes the entire monotone polar of a class of graphs, gives a necessary and sufficient condition for their maximal monotonicity, and shows how a positive rank-one perturbation produces a nonmaximal sum under this condition. By coupling triangular operators with a Lipschitz curve, we verify the theorem's hypotheses and its maximality criterion on $c_0$, thereby obtaining a counterexample to the conjecture. We then construct a bounded linear surjection from $\ell^1$ onto $c_0$ and use it to obtain the counterexample on $\ell^1$.





\end{abstract}


\noindent\textbf{Keywords:} maximally monotone operator; sum theorem;
nonreflexive Banach space; monotone polar; rank-one operator.

\section{Introduction}
\label{sec:introduction}
Let $X$ be a real Banach space, let $X^*$ be its continuous dual, and let $A,B:X\rightrightarrows X^*$ be maximally monotone operators. Their pointwise sum is defined by $(A+B)x=\{a+b:a\in Ax,\ b\in Bx\}$. Rockafellar's sum theorem \cite{rockafellar1970} states that $A+B$ is maximally monotone when $X$ is reflexive and
\begin{equation}
\dom A\cap\operatorname{int}(\dom B)\ne\varnothing.
\label{eq:original-cq}
\end{equation}
The unrestricted sum question asks whether Eq. (\ref{eq:original-cq}) alone is sufficient on an arbitrary real Banach space. The interior here is taken in the norm topology of $X$, and maximality is taken in the full pair $X\times X^*$. Bo\c{t}, Bueno, and Simons \cite{botbuenosimons2016} show that the unrestricted sum question is equivalent to its restriction to pairs with bounded common domain.

In general Banach spaces, sum theorems have been established under additional assumptions on the operators. Yao \cite{yao2010} proves maximality when $A$ is of type~(FPV) and $B$ has full domain. Borwein and Yao \cite{borweinyao2012} prove maximality under Eq. (\ref{eq:original-cq}) when $A$ is a linear relation. Voisei \cite{voisei2010} treats the sum of an operator of type~(FPV) with convex domain and a normal-cone operator. Verona and Verona \cite[Main Result~1.1(b)]{veronaverona2022} prove the sum theorem under Eq. (\ref{eq:original-cq}) when $A$ is $\mathcal C_0$-maximal: $A+N_C$ is maximally monotone for every closed convex set $C$ with $\dom A\cap\operatorname{int}C\ne\varnothing$, where $N_C$ denotes the normal-cone operator of $C$. Bauschke, Borwein, Wang, and Yao \cite{bauschkeetal2011} develop triangular operators and nonlinear examples. We use the triangular map of Bauschke et al. as a building block for our construction on $c_0$.

Sufficient conditions for maximality of sums have also been obtained through convex representations of monotone operators. Fitzpatrick \cite{fitzpatrick1988} associates a convex function with a monotone graph, and Burachik and Svaiter \cite{burachiksvaiter2002} study the associated family of convex representatives through enlargements. Using convex representations, Penot \cite{penot2004} develops results for compositions and sums in reflexive spaces, and Marques Alves and Svaiter \cite{marquessvaiter2012} prove a sum theorem in general Banach spaces under a qualification on the domains of Fitzpatrick representations.



In this paper, we construct counterexamples to Rockafellar's sum conjecture on $c_0$ and standard $\ell^1$. Our general construction theorem computes the entire monotone polar of a class of graphs, characterizes their maximality, and identifies a positive rank-one perturbation whose addition makes the sum nonmaximal. To satisfy the theorem's assumptions on $c_0$, we couple copies of the triangular map through a Lipschitz curve, and then obtain a counterexample to the conjecture. A bounded linear surjection then transfers this counterexample to standard $\ell^1$.

\subsection{Contributions}
The contributions of this paper are as follows.

(1) We establish a general construction theorem for counterexamples to Rockafellar's sum conjecture (Theorem~\ref{thm:endpoint-criterion}). Under Assumption~\ref{ass:construction}, the theorem computes the entire monotone polar of the graph in Eq. (\ref{eq:abs-full-graph}) and characterizes its maximality by the criterion in Eq. (\ref{eq:abs-endpoint-exclusion}). When this criterion holds, the graph defines a maximally monotone operator whose sum with a specified everywhere-defined positive rank-one operator is not maximally monotone, although the two operators satisfy Eq. (\ref{eq:original-cq}).

(2) We construct an explicit counterexample on $c_0$ (Example~\ref{thm:C}) by verifying the assumptions of this theorem. The second operator is bounded, positive, rank-one, and defined everywhere, and the first operator has nonempty domain. Thus the original interior-domain condition holds. We also establish a finite radial bound for the first operator.

(3) We construct a counterexample on standard $\ell^1$ (Example~\ref{thm:L}) by pulling back the first operator through a specified bounded linear surjection onto $c_0$. The pullback lemma proves maximality in $\ell^1\times\ell^\infty$, and the adjoint identity $\pair{u}{Q^*a}=\pair{Qu}{a}$ transfers the inequalities establishing nonmaximality of the sum.

The paper is organized as follows. Section~\ref{sec:definitions} introduces the notation and basic definitions. Section~\ref{sec:endpoint-criterion} proves the structural theorem and the pullback lemma. Section~\ref{sec:counterexamples} constructs the counterexamples on $c_0$ and standard $\ell^1$ and derives their stated consequences. Section~\ref{sec:conclusion} concludes the paper.


\section{Preliminaries}
\label{sec:definitions}
All scalars are real. Put $\N=\{1,2,\ldots\}$ and $\Nzero=\{0,1,\ldots\}$. For a Banach space $X$, we denote by $X^*$ its continuous dual, the space of continuous linear functionals on $X$. We write $\pair{x}{a}=a(x)$ for $x\in X$ and $a\in X^*$. For a real Hilbert space $H$, its inner product is denoted by $(u,v)_H$. We use the norm topology for interiors and closures unless stated otherwise.

For a bounded linear map $P:X\to X^*$, we use positive to mean $\pair{x}{Px}\ge0$ for every $x\in X$, and rank one to mean $\dim(\operatorname{ran}P)=1$.

\begin{definition}[Monotonicity and the monotone polar]
For a graph $G\subset X\times X^*$, define
\begin{equation}
G^\mu=\{(z,p)\in X\times X^*:
 \pair{z-x}{p-a}\ge0\quad\forall(x,a)\in G\}.
\label{eq:D1}
\end{equation}
The graph is monotone when $\pair{x-y}{a-b}\ge0$ for all $(x,a),(y,b)\in G$, and maximally monotone when it has no proper monotone extension in this same full dual pair.
\end{definition}

\begin{proposition}
A monotone graph $G\subset X\times X^*$ is maximally monotone if and only if $G=G^\mu$.
\end{proposition}
\begin{proof}
A point in $G^\mu\setminus G$ can be adjoined to $G$, and every point of a monotone extension belongs to $G^\mu$.
\end{proof}

We write $\dom G=\{x:\exists a,\ (x,a)\in G\}$ and $\operatorname{ran}G=\{a:\exists x,\ (x,a)\in G\}$.


\begin{definition}[Radial bound at the origin]
For a nonempty graph $G$ with $0\notin\dom G$, define
\begin{equation}
\mathcal V(G)=\sup_{(x,a)\in G}
 \frac{\max\{0,-\pair{x}{a}\}}{\|x\|}.
\label{eq:D3}
\end{equation}
\end{definition}
For $G=\gra T$ as in the preceding definition, the quantity of Verona and Verona \cite[p.~1004]{verona2008} satisfies
\[
\begin{aligned}
L(0,0,T)
&=\max\left\{0,\sup_{(x,a)\in G}\frac{-\pair{x}{a}}{\|x\|}\right\}\\
&=\sup_{(x,a)\in G}\frac{\max\{0,-\pair{x}{a}\}}{\|x\|}
=\mathcal V(G).
\end{aligned}
\]
For every $C\ge0$, Eq. (\ref{eq:D3}) gives
\[
\mathcal V(G)\le C
\quad\Longleftrightarrow\quad
\pair{x}{a}\ge-C\|x\|\quad\text{for all }(x,a)\in G.
\]
Thus, when finite, $\mathcal V(G)$ is the least such constant. The bounds for our two operators are proved in Examples~\ref{thm:C} and~\ref{thm:L}, using Eqs. (\ref{eq:C9}) and (\ref{eq:L6}), respectively.

\section{Maximal monotonicity and nonmaximal sums}
\label{sec:endpoint-criterion}

In this section, we establish a general construction theorem for maximally monotone operators whose sums with an everywhere-defined positive rank-one operator are not maximally monotone. We also prove a lemma that transfers maximal monotonicity through a bounded linear surjection.

% Let $X$ be a real Banach space and let $H_0$ be a real Hilbert space.
% We consider graphs of the form $G=\{(La,a):a\in D\}$, where
% $L:X^*\to X$ is bounded and linear, and $D\subset X^*$ is specified
% by the constraints below.




Let $X$ be a real Banach space.
\begin{definition}[Hilbert spaces and norms]
\label{def:hilbert-spaces-norms}
Let $H_0$ be a real Hilbert space with inner product $(\cdot,\cdot)_{H_0}$ and induced norm
\[
\|u\|_{H_0}=\sqrt{(u,u)_{H_0}}\qquad(u\in H_0).
\]
Define $H=\R\oplus H_0$ to be $\R\times H_0$ equipped with the inner product
\[
((s,u),(t,v))_H=st+(u,v)_{H_0}
\qquad(s,t\in\R,\ u,v\in H_0).
\]
Its induced norm is
\[
\|(s,u)\|_H=\sqrt{s^2+\|u\|_{H_0}^2}.
\]
\end{definition}
Next, we give the following definition.

\begin{definition}[The linear maps]
\label{def:construction-maps}
Let $E:X^*\to X$ and $M:X^*\to H$ be bounded linear maps, and let $g\in X^*\setminus\{0\}$. Define
\[
\begin{aligned}
h&=Eg,\\
r(a)&=\pair{h}{a},\\
La&=-Ea+r(a)h,\\
K&=\ker M\cap\ker r.
\end{aligned}
\]
\end{definition}

We use $r(a)$ as a scalar parameter and constrain $Ma$ to lie on a prescribed curve at that parameter. This specifies the graph as follows.

\begin{definition}[The constrained graph]
\label{def:construction-graph}
Let $\omega:(0,\infty)\to H_0$ be $\ell$-Lipschitz, where $0\le\ell\le1$. Define
\[
\begin{aligned}
J&=(-\infty,-1)\cup(1,\infty),\\
F(t)&=(\sgn t,\omega(|t|-1))\qquad(t\in J), 
\end{aligned}
\]
where \[
\operatorname{sgn}(t)=
\begin{cases}
1, & t>0,\\
0, & t=0,\\
-1, & t<0.
\end{cases}
\]. 
The admissible dual variables and the graph are
\begin{equation}
\begin{aligned}
D&=\{a\in X^*:r(a)\in J,\ Ma=F(r(a))\},\\
G&=\{(La,a):a\in D\}.
\end{aligned}
\label{eq:abs-full-graph}
\end{equation}
\end{definition}

To compute the entire monotone polar of $G$, we impose the following conditions on the linear maps, the curve, and its realizability in $X^*$.
\begin{assumption}\label{ass:construction}
For the maps and graph in Definitions~\ref{def:construction-maps} and~\ref{def:construction-graph}, assume:
\begin{enumerate}
\renewcommand{\theenumi}{H\arabic{enumi}}
\renewcommand{\labelenumi}{(\theenumi)}
\item\label{ass:bilinear} $Mg=0$, and
\begin{equation}
\pair{Ea}{b}+\pair{Eb}{a}=2(Ma,Mb)_H
\qquad(a,b\in X^*).
\label{eq:abs-bilinear}
\end{equation}

\item\label{ass:annihilator} The subspace $K$ satisfies
\begin{equation}
\{x\in X:\pair{x}{k}=0\ \text{for all }k\in K\}=\R h.
\label{eq:abs-annihilator}
\end{equation}

\item\label{ass:curve-limit} There exists $\eta\in H_0$ such that
\begin{equation}
\begin{gathered}
\omega(s)\longrightarrow\eta\text{ in }H_0
\quad\text{as }s\downarrow0,\\
0<S:=\|\eta\|_{H_0}^2<1,\\
\|\omega(s)\|_{H_0}^2\le S\qquad(s>0).
\end{gathered}
\label{eq:abs-tail}
\end{equation}

\item\label{ass:realization} For every $t\in J$, there exists $a^t\in X^*$
satisfying $r(a^t)=t$ and $Ma^t=F(t)$.
\end{enumerate}
\end{assumption}
By Assumption~\ref{ass:construction}\textup{(\ref{ass:realization})}, each parameter fibre $\{a\in X^*:r(a)=t,\ Ma=F(t)\}$ is nonempty and equals $a^t+K$. Definition~\ref{def:construction-graph} includes every point in each fibre. Combining Eq. (\ref{eq:abs-bilinear}) and the Lipschitz bound on $\omega$ with these affine fibres, we obtain an explicit formula for $G^\mu$. The following theorem shows that, under Assumption~\ref{ass:construction}, $G$ is maximally monotone if and only if no $(p,\tau)\in X^*\times[-1,1]$ satisfies $r(p)=\tau$ and $Mp=(\tau,\eta)$. When this condition holds, adding the operator $\mathcal P x=\pair{x}{g}g$ to the operator with graph $G$ produces a nonmaximal sum. Therefore, it gives a general construction criterion for counterexamples satisfying the interior-domain condition in Eq. (\ref{eq:original-cq}) as follows.
% The theorem below determines
% when such a graph is maximally monotone and when adding a positive
% rank-one operator produces a nonmaximal sum. 

\begin{theorem}[Construction of nonmaximal sums]
\label{thm:endpoint-criterion}
For the graph $G$ in Definition~\ref{def:construction-graph}, suppose Assumption~\ref{ass:construction} holds and define
\[
\widehat F(t)=
\begin{cases}
F(t),&|t|>1,\\
(t,\eta),&|t|\le1.
\end{cases}
\]
Then the following statements hold.
\begin{enumerate}
\renewcommand{\theenumi}{\roman{enumi}}
\renewcommand{\labelenumi}{(\theenumi)}
\item\label{thm:polar-part} The entire monotone polar of $G$ in $X\times X^*$ is
\begin{equation}
G^\mu=\{(Lp,p):p\in X^*,\ Mp=\widehat F(r(p))\}.
\label{eq:abs-entire-polar}
\end{equation}
This graph is the unique maximally monotone extension of $G$.

\item\label{thm:maximality-part} The graph $G$ is maximally monotone if and only if
\begin{equation}
\nexists(p,\tau)\in X^*\times[-1,1]:
\quad r(p)=\tau,\quad Mp=(\tau,\eta).
\label{eq:abs-endpoint-exclusion}
\end{equation}

\item\label{thm:sum-part} Suppose Eq. (\ref{eq:abs-endpoint-exclusion}) holds.
Let $\mathcal A$ be the operator with graph $G$, and define $\mathcal P x=\pair{x}{g}g$. Then
\[
0\notin\dom\mathcal A,\qquad \mathcal V(G)\le S\|g\|.
\]
Both $\mathcal A$ and $\mathcal P$ are maximally monotone and satisfy
\[
\dom\mathcal A\cap\operatorname{int}(\dom\mathcal P)\ne\varnothing.
\]
Their sum is not maximally monotone. More precisely,
\[
(0,0)\in
\bigl(\gra(\mathcal A+\mathcal P)\bigr)^\mu
\setminus\gra(\mathcal A+\mathcal P).
\]
\end{enumerate}
\end{theorem}
\begin{proof}
\noindent\textit{\textup{(\ref{thm:polar-part})} The monotone polar and its maximality.}
Eq. (\ref{eq:abs-bilinear}) at $a=b=g$ gives $r(g)=0$. Using $Mg=0$ again gives
\begin{equation}
\begin{aligned}
\pair{Ea}{g}&=-r(a),\\
\pair{La}{g}&=r(a),\\
\pair{La-Lb}{a-b}&=|r(a)-r(b)|^2-\|Ma-Mb\|^2.
\end{aligned}
\label{eq:abs-pairing}
\end{equation}
Taking $t=2$ in Assumption~\ref{ass:construction}\textup{(\ref{ass:realization})} gives $r(a^2)=2$, so $h\ne0$.

Extend $\omega$ to zero by $\omega(0)=\eta$. Its Lipschitz bound persists by the norm limit. Define $\chi(t)=\max\{-1,\min\{1,t\}\}$ and $\zeta(t)=\max\{|t|-1,0\}$. Then $\widehat F(t)=(\chi(t),\omega(\zeta(t)))$. On each of $[0,\infty)$ and $(-\infty,0]$, the definitions give
\[
|\chi(t)-\chi(u)|+|\zeta(t)-\zeta(u)|=|t-u|.
\]
For $t\ge0\ge u$, they give
\[
\begin{aligned}
|\chi(t)-\chi(u)|+|\zeta(t)-\zeta(u)|
&\le \chi(t)-\chi(u)+\zeta(t)+\zeta(u)\\
&=t-u.
\end{aligned}
\]
Interchanging $t,u$ covers the remaining case. Therefore
\[
|\chi(t)-\chi(u)|+|\zeta(t)-\zeta(u)|\le|t-u|,
\]
and
\[
\|\widehat F(t)-\widehat F(u)\|^2
\le|\chi(t)-\chi(u)|^2+\ell^2|\zeta(t)-\zeta(u)|^2
\le|t-u|^2.
\]
Eq. (\ref{eq:abs-pairing}) proves monotonicity of both $G$ and the set on the right of Eq. (\ref{eq:abs-entire-polar}). It also proves that this latter set is contained in $G^\mu$.

Conversely, let $(x,p)\in G^\mu$. For $a\in D$ and $t=r(a)$, Eq. (\ref{eq:abs-bilinear}) gives
\begin{align*}
\pair{x-La}{p-a}
={}&\pair{x}{p}-\pair{x+Ep}{a}
       +2(Ma,Mp)_H-t r(p)+t^2-\|Ma\|^2.
\end{align*}
Replacing $a$ by $a+\theta k$, for all $\theta\in\R$ and $k\in K$, forces $\pair{x+Ep}{k}=0$. Thus, $x=-Ep+v h$ by Eq. (\ref{eq:abs-annihilator}). With
\[
\tau=\frac{v+r(p)}2,\qquad \nu=\frac{v-r(p)}2,
\]
the same expansion reduces exactly to
\begin{equation}
\pair{x-La}{p-a}=(t-\tau)^2-\nu^2-\|F(t)-Mp\|^2.
\label{eq:abs-polar-square}
\end{equation}
By Assumption~\ref{ass:construction}\textup{(\ref{ass:realization})}, the polar inequality holds for every $t\in J$. We distinguish two cases.

\emph{Case 1: $|\tau|>1$.} Taking $t=\tau$ in Eq. (\ref{eq:abs-polar-square}) gives $\nu=0$ and $Mp=F(\tau)$. Hence $v=r(p)=\tau$ and $x=Lp$.

\emph{Case 2: $|\tau|\le1$.} Write $Mp=(b,z)\in H$. Letting $t\to1^+$ and $t\to-1^-$, respectively, in Eq. (\ref{eq:abs-polar-square}), and using Assumption~\ref{ass:construction}\textup{(\ref{ass:curve-limit})}, gives
\[
\begin{aligned}
(b-1)^2+\|z-\eta\|^2+\nu^2&\le(1-\tau)^2,\\
(b+1)^2+\|z-\eta\|^2+\nu^2&\le(1+\tau)^2.
\end{aligned}
\]
Multiplying the two inequalities by the nonnegative weights $(1+\tau)/2$ and $(1-\tau)/2$, respectively, and adding gives
\[
(b-\tau)^2+\|z-\eta\|^2+\nu^2\le0.
\]
Consequently $Mp=(\tau,\eta)$, $r(p)=v=\tau$, and $x=Lp$. This proves Eq. (\ref{eq:abs-entire-polar}). Here $p\in X^*$ is fixed; only $F(t)$ is passed to the limit in $H$.

We have proved that $G^\mu$ is monotone. Any point related to it is related to $G$ and therefore belongs to $G^\mu$, so $G^\mu$ is maximally monotone. Since every monotone extension of $G$ lies in $G^\mu$, a maximal one must equal it.

\noindent\textit{\textup{(\ref{thm:maximality-part})} Maximality of $G$.}
By part~\textup{(\ref{thm:polar-part})}, $G^\mu\setminus G$ consists exactly of the points $(Lp,p)$ for which $r(p)\in[-1,1]$ and $Mp=(r(p),\eta)$, as in Eq. (\ref{eq:abs-endpoint-exclusion}). Thus $G=G^\mu$ if and only if these equations have no solution.

\noindent\textit{\textup{(\ref{thm:sum-part})} The radial bound and nonmaximality of the sum.}
Assume Eq. (\ref{eq:abs-endpoint-exclusion}). For each $a\in D$,
\[
\begin{aligned}
|\pair{La}{g}|&=|r(a)|>1,\\
\|La\|&>1/\|g\|,\\
\pair{La}{a}&=r(a)^2-1-\|\omega(|r(a)|-1)\|^2>-S.
\end{aligned}
\]
Thus $0\notin\dom\mathcal A$ and, for every $a\in D$,
\[
\frac{\max\{0,-\pair{La}{a}\}}{\|La\|}
\le\frac{S}{\|La\|}<S\|g\|.
\]
Taking the supremum proves $\mathcal V(G)\le S\|g\|$. The map $\mathcal P$ is bounded, linear, positive, nonzero, and defined on all of $X$. It is also the subdifferential of the continuous convex function $x\mapsto\pair{x}{g}^2/2$, in accordance with \cite{rockafellar1970subdiff}. We prove its maximality directly. Its monotonicity follows from $\pair{x-y}{\mathcal Px-\mathcal Py}=\pair{x-y}{g}^2$. For an arbitrary point $(z,z^*)$ in its polar, tests at $z+tv$ give
\[
-t\pair{v}{z^*-\mathcal Pz}+t^2\pair{v}{g}^2\ge0
\qquad(t\in\R,\ v\in X).
\]
Both signs of arbitrarily small $t$ force $z^*=\mathcal Pz$. Thus $\mathcal P$ is maximally monotone, and $La^2\in\dom\mathcal A$ proves $\dom\mathcal A\cap\operatorname{int}\dom\mathcal P\ne\varnothing$. For every $a\in D$,
\[
\pair{La}{a+\mathcal PLa}
=2r(a)^2-1-\|\omega(|r(a)|-1)\|^2>1-S>0.
\]
The sum is monotone and its domain does not contain zero. Adjoining $(0,0)$ is a proper monotone extension.
\end{proof}

Theorem~\ref{thm:endpoint-criterion} provides a criterion for constructing the first operator together with a rank-one perturbation that makes its sum nonmaximal. For a graph satisfying Assumption~\ref{ass:construction}, part~\textup{(\ref{thm:polar-part})} identifies every possible additional polar point. Part~\textup{(\ref{thm:maximality-part})} therefore reduces maximality to excluding $p\in X^*$ and $\tau\in[-1,1]$ solving
\[
\begin{aligned}
r(p)&=\tau,\\
Mp&=(\tau,\eta).
\end{aligned}
\]
Once this is proved, part~\textup{(\ref{thm:sum-part})} supplies the second operator $\mathcal P$ and the point $(0,0)$ witnessing nonmaximality of the sum. Thus, to obtain the example on $c_0$, we will choose the maps and the curve so that the system has no solution in the continuous dual of $c_0$. The following lemma then transfers that example to standard $\ell^1$.

\begin{lemma}[Pullback by a bounded surjection]
\label{lem:surjection-transport}
Let $U,V$ be real Banach spaces and let $Q:U\to V$ be a bounded linear surjection. Suppose $C_Q>0$ and every $y\in V$ has a preimage $u\in U$ with $\|u\|\le C_Q\|y\|$. If $\mathcal M:V\rightrightarrows V^*$ is maximally monotone, then
\[
\gra(Q^*\mathcal M Q)
=\{(u,Q^*a):(Qu,a)\in\gra\mathcal M\}
\]
is maximally monotone in the full pair $U\times U^*$, where $Q^*:V^*\to U^*$ is the adjoint of $Q$, defined by $\pair{u}{Q^*a}=\pair{Qu}{a}$ for $u\in U$ and $a\in V^*$.
\end{lemma}
\begin{proof}
A maximally monotone graph is nonempty, since the empty graph admits a one-point monotone extension. Surjectivity therefore makes the pullback graph nonempty. For two of its points,
\[
\pair{u-v}{Q^*(a-b)}=\pair{Qu-Qv}{a-b}\ge0,
\]
so it is monotone.

Let $(v,v^*)\in U\times U^*$ be related to this entire graph. Fix $(u,Q^*a)$ in the graph. For each $k\in\ker Q$, all $(u+t k,Q^*a)$ belong to it, and their polar inequalities have the form
\[
\pair{v-u}{v^*-Q^*a}-t\pair{k}{v^*}\ge0
\qquad(t\in\R).
\]
Thus $v^*$ annihilates $\ker Q$. Define $p(y)=\pair{u}{v^*}$ when $Qu=y$. Differences of preimages prove well-definedness, and linear combinations of preimages prove linearity. The prescribed lifting bound gives $|p(y)|\le C_Q\|v^*\|\|y\|$, so $p\in V^*$ and $v^*=Q^*p$. For every $(y,a)\in\gra\mathcal M$, choose $u\in U$ with $Qu=y$. Then $(u,Q^*a)$ belongs to the pullback graph. The remaining inequalities are
\[
\pair{Qv-y}{p-a}\ge0\qquad((y,a)\in\gra\mathcal M).
\]
Maximality of $\mathcal M$ gives $(Qv,p)\in\gra\mathcal M$, and $(v,v^*)$ belongs to the pullback graph.
\end{proof}

\section{Counterexamples on $c_0$ and standard $\ell^1$}
\label{sec:counterexamples}
In this section, we apply Theorem~\ref{thm:endpoint-criterion} to construct a counterexample on $c_0$ and then use Lemma~\ref{lem:surjection-transport} to obtain a counterexample on standard $\ell^1$.


% We first verify the structural theorem for explicit operators on $c_0$.
% A bounded surjection then defines operators on standard $\ell^1$, whose
% maximality and sum failure are proved in the full continuous dual pair.

\subsection{A counterexample on $c_0$}

Put $I=\Nzero\times\N$ and $Y=c_0(I)$ with the supremum norm. For every $x\in Y$ and $\epsilon>0$, only finitely many $(b,j)\in I$ satisfy $|x_{b,j}|\ge\epsilon$. Its continuous dual is $\ell^1(I)$, with $\pair{x}{a}=\sum_{b,j}x_{b,j}a_{b,j}$. We also use this pairing for $x\in\ell^\infty(I)$ and $a\in\ell^1(I)$. The unit coordinate vectors are denoted by $e_{b,j}$.

We index coordinates by blocks. The map $m$ records the sum in each block, and $E$ applies a triangular operator within that block. Its single-block formula is the $\alpha=(1,1,\ldots)$ case of \cite[author version, Example~4.1]{bauschkeetal2011}; the additional constraint below couples the block sums through $F$.
\begin{definition}[Block maps]
For $a\in\ell^1(I)$, define
\begin{equation}
\begin{aligned}
(ma)_b&=\sum_{j\ge1}a_{b,j},\\
(Ea)_{b,j}&=a_{b,j}+2\sum_{k>j}a_{b,k}.
\end{aligned}
\label{eq:D4}
\end{equation}
Here $m:\ell^1(I)\to\ell^1(\Nzero)\subset\ell^2(\Nzero)$ and $E:\ell^1(I)\to Y$. Put
\begin{equation}
\begin{aligned}
g&=e_{0,1}-e_{0,2}\in Y^*,\\
h&=Eg=-e_{0,1}-e_{0,2}\in Y,\\
r(a)&=\pair{h}{a}=-a_{0,1}-a_{0,2},\\
La&=-Ea+r(a)h.
\end{aligned}
\label{eq:D5}
\end{equation}
\end{definition}

\begin{definition}[The curve of block sums]
For $s>0$ and $n\ge1$, write
\begin{equation}
\begin{aligned}
\omega_s(n)&=\frac{\max\{1-sn^{1/4},0\}}{4n},\\
W(s)&=\sum_{n\ge1}\omega_s(n)^2,\\
\sigma&=\sum_{n\ge1}\frac1{16n^2}=\frac{\pi^2}{96}.
\end{aligned}
\label{eq:D6}
\end{equation}
Each $\omega_s$, $s>0$, has finite support. Set $\omega_0=(1/(4n))_{n\ge1}\in\ell^2\setminus\ell^1$. Define
\begin{equation}
\begin{aligned}
J&=(-\infty,-1)\cup(1,\infty),\\
F(t)&=(\sgn t,\omega_{|t|-1})\in\ell^1(\Nzero)\quad(t\in J),\\
D&=\{a\in\ell^1(I):r(a)\in J,\ ma=F(r(a))\}.
\end{aligned}
\label{eq:D7}
\end{equation}

\end{definition}

\begin{definition}[The two operators on $c_0(I)$]
The operators and kernel are
\begin{equation}
\begin{aligned}
\gra A&=\{(La,a):a\in D\},\\
K&=\{k\in\ell^1(I):mk=0,\ r(k)=0\},\\
Px&=\pair{x}{g}g\quad(x\in Y).
\end{aligned}
\label{eq:D8}
\end{equation}
For each $t\in J$, a vector satisfying $r(a^t)=t$ and $ma^t=F(t)$ is
\begin{equation}
a^t=-t e_{0,1}+(t+\sgn t)e_{0,3}
       +\sum_{n\ge1}\omega_{|t|-1}(n)e_{n,1}.
\label{eq:D9}
\end{equation}
\end{definition}

% Coordinate verification follows the definitions in the first example.
\label{sec:supporting}
We now verify the assumptions of Theorem~\ref{thm:endpoint-criterion} for the maps in Eqs. (\ref{eq:D4})--(\ref{eq:D9}). The first three lemmas supply the pairing identity, the exact primal annihilator and the complete parameter fibres.

\begin{lemma}[The coordinate pairing identity]
\label{lem:coordinate-pairing}
The maps $m:\ell^1(I)\to\ell^1(\Nzero)\subset\ell^2(\Nzero)$ and $E,L:\ell^1(I)\to c_0(I)$ are bounded. For all $a,c\in\ell^1(I)$, Eqs. (\ref{eq:C1}) and (\ref{eq:C2}) below hold.
\end{lemma}
\begin{proof}
\label{sec:C1}
For $a\in\ell^1(I)$, $\|ma\|_1\le\|a\|_1$ and $\|Ea\|_\infty\le2\|a\|_1$. If $a$ has finite support, then $Ea$ has finite support: only finitely many blocks occur, and within a block no output occurs after the last nonzero input coordinate. Approximation by finite-support vectors and the bound prove $Ea\in c_0(I)$. Thus $L$ also maps $\ell^1(I)$ boundedly into $c_0(I)$.

For $a,c\in\ell^1(I)$, absolute convergence permits exchanging the double sums. In each block the diagonal terms and the two triangular off-diagonal sums give
\begin{equation}
\pair{Ea}{c}+\pair{Ec}{a}
 =2\sum_b\Big(\sum_j a_{b,j}\Big)\Big(\sum_j c_{b,j}\Big)
 =2(ma,mc)_{\ell^2}.
\label{eq:C1}
\end{equation}
The absolute sum of all products is at most $\|a\|_1\|c\|_1$ before the factor $2$, so the identity holds for all $a,c\in\ell^1(I)$. Since $mg=0$, $h=Eg$, and $r(g)=0$, Eq. (\ref{eq:C1}) implies $\pair{Ea}{g}=-r(a)$. Consequently
\begin{equation}
\begin{aligned}
\pair{La}{a}&=r(a)^2-\|ma\|_2^2,\\
\pair{La}{g}&=r(a),\\
\pair{La-Lc}{a-c}&=(r(a)-r(c))^2-\|ma-mc\|_2^2.
\end{aligned}
\label{eq:C2}
\end{equation}
\end{proof}

\begin{lemma}[The annihilator of $K$ in $c_0(I)$]
\label{lem:coordinate-annihilator}
For $K$ in Eq. (\ref{eq:D8}),
\[
\{z\in c_0(I):\pair{z}{k}=0\ \text{for every }k\in K\}=\R h.
\]
\end{lemma}
\begin{proof}
Let $z\in c_0(I)$ annihilate $K$. In a block $b\ge1$, every difference $e_{b,j}-e_{b,k}$ is in $K$. Thus $z$ is constant in that entire block. Since $z\in c_0(I)$, this constant must be zero. In block zero, differences between indices $j,k\ge3$ belong to $K$, so the common tail value is also zero. Finally, $g=e_{0,1}-e_{0,2}$ belongs to $K$, forcing $z_{0,1}=z_{0,2}$. It follows that $z$ is a scalar multiple of $h$. Conversely, $\pair{h}{k}=r(k)=0$ proves that every multiple of $h$ annihilates $K$.
\end{proof}

\begin{lemma}[Properties of the curve and its parameter fibres]
\label{lem:coordinate-profile}
The map $s\mapsto\omega_s$ from $(0,\infty)$ to $\ell^2(\N)$ satisfies Eq. (\ref{eq:abs-tail}) with $\eta=\omega_0$, $S=\sigma<1/8$ and a Lipschitz constant strictly less than one. For every $t\in J$, the vector $a^t$ in Eq. (\ref{eq:D9}) belongs to $D$, and its entire parameter fibre is $a^t+K$. No $p\in\ell^1(I)$ and $\tau\in[-1,1]$ satisfy $r(p)=\tau$ and $mp=(\tau,\omega_0)$.
\end{lemma}
\begin{proof}
\label{sec:C2}
For each $s>0$, $\omega_s(n)=0$ whenever $n\ge s^{-4}$. Also $0\le\omega_s(n)\le1/(4n)$. Thus, $W(s)\le\sigma<1/8$, and dominated convergence in $\ell^2$ gives $\omega_s\to\omega_0$ and $W(s)\to\sigma$ as $s$ decreases to zero. The comparison vector $\omega_0$ is not in $\ell^1$ by divergence of the harmonic series.

Since $v\mapsto\max\{v,0\}$ is $1$-Lipschitz on $\R$,
\begin{equation}
\|\omega_s-\omega_q\|_2^2\le c|s-q|^2,\qquad
c=\frac1{16}\sum_{n\ge1}n^{-3/2}<\frac3{16}<1.
\label{eq:C3}
\end{equation}
The strict numerical bound follows from $\sum_{n\ge1}n^{-3/2}<1+\int_1^\infty x^{-3/2}\,dx=3$. For $t,u$ on the same branch of $J$, Eq. (\ref{eq:C3}) gives $\|F(t)-F(u)\|_2\le|t-u|$. For $t=1+s$ and $u=-1-q$, the squared distance is at most $4+c(s-q)^2\le(2+s+q)^2$, and the opposite orientation is identical. Thus $F$ is $1$-Lipschitz on all of $J$.

The finitely supported vector $a^t$ in Eq. (\ref{eq:D9}) satisfies $r(a^t)=t$ and $ma^t=F(t)$, proving nonemptiness for every $t\in J$. The complete fibre $\{a:r(a)=t,\ ma=F(t)\}$ is precisely $a^t+K$. For every $p\in\ell^1(I)$, the vector $mp$ is in $\ell^1(\Nzero)$. Since $\omega_0\notin\ell^1(\N)$, the last assertion follows.
\end{proof}

% The first counterexample follows its coordinate lemmas.
\label{sec:c0}
\label{sec:statements}
\begin{example}[A counterexample on $c_0$]\label{thm:C}
Let $I=\Nzero\times\N$ and $Y=c_0(I)$ with its usual supremum norm and full dual $Y^*=\ell^1(I)$. Let $A:Y\rightrightarrows Y^*$ and the bounded positive rank-one map $P:Y\to Y^*$ be defined by Eq. (\ref{eq:D8}). Then:
\begin{enumerate}
\item $A$ and $P$ are maximally monotone in $Y\times Y^*$.
\item $\dom P=Y$ and $\dom A\ne\varnothing$, so
$\dom A\cap\operatorname{int}(\dom P)\ne\varnothing$.
\item Every $(x,a)\in\gra A$ satisfies $\|x\|_\infty>1/2$ and
$\pair{x}{a}\ge-2\sigma\|x\|_\infty$, where $\sigma=\pi^2/96<1/8$.
\item Every $(x,b)\in\gra(A+P)$ satisfies
$\pair{x}{b}>1-\sigma>7/8$.
\item $(0,0)$ is monotonically related to the entire graph of $A+P$ and
does not belong to that graph. Thus, the $A+P$ is not maximally monotone.
\end{enumerate}
\end{example}
\begin{proof}
Lemmas~\ref{lem:coordinate-pairing}--\ref{lem:coordinate-profile} verify every assumption of Theorem~\ref{thm:endpoint-criterion} with $X=Y$, $M=m$, $H_0=\ell^2(\N)$, $\omega(s)=\omega_s$, $\eta=\omega_0$ and $S=\sigma$. In particular, $g\ne0$ and $mg=0$ by Eq. (\ref{eq:D5}), all parameter fibres are present and the exact primal annihilator is $\R h$.

By Eq. (\ref{eq:C2}), the Lipschitz constraint proves monotonicity for every $a,c\in D$. The formula for the monotone polar in Eq. (\ref{eq:abs-entire-polar}) and the nonsummability of $\omega_0$ then give $(\gra A)^\mu=\gra A$: no $p\in\ell^1(I)$ can satisfy $mp=(\tau,\omega_0)$ for $|\tau|\le1$. This proves maximality in $c_0(I)\times\ell^1(I)$.

\label{sec:C5}
For every $a\in D$, $t=r(a)$ satisfies $|t|>1$. By Eq. (\ref{eq:C2}) and $\|g\|_1=2$,
\begin{equation}
\begin{aligned}
\|La\|_\infty&\ge|t|/2>1/2,\\
\pair{La}{a}&=t^2-1-W(|t|-1)>-\sigma.
\end{aligned}
\label{eq:C9}
\end{equation}
Eq. (\ref{eq:C9}) gives $A(0)=\varnothing$ and, for every $a\in D$,
\[
\frac{\max\{0,-\pair{La}{a}\}}{\|La\|_\infty}
\le \frac{\sigma}{\|La\|_\infty}
<2\sigma.
\]
Hence $\pair{La}{a}\ge-2\sigma\|La\|_\infty$ for every $a\in D$, and taking the supremum gives $\mathcal V(\gra A)\le2\sigma<1/4$.

\label{sec:C6}
The rank-one assertion and maximality of $P$ follow from the direct proof in Theorem~\ref{thm:endpoint-criterion}. Here $\|g\|_1=2$, so $\|P\|\le4$, and $Pe_{0,1}=g\ne0$. The vector $a^2=-2e_{0,1}+3e_{0,3}$ gives $La^2=-6e_{0,1}-8e_{0,2}-3e_{0,3}$. Thus $\dom A\cap\operatorname{int}(\dom P)\ne\varnothing$, with $\operatorname{int}(\dom P)=Y$. Every sum point has the form $(La,a+tg)$, $a\in D$ and $t=r(a)$, and
\begin{equation}
\pair{-La}{-a-tg}=2t^2-1-W(|t|-1)>1-\sigma>7/8.
\label{eq:C11}
\end{equation}
This proves that $(0,0)$ is related to every sum point. The sum is monotone as the sum of two monotone operators; $0\notin\dom(A+P)$ by Eq. (\ref{eq:C9}). Adding $(0,0)$ is therefore a proper monotone extension.

\end{proof}

% Consequences of the first counterexample.
\label{sec:domain-comparison}

% \begin{corollary}\label{prop:domain-gap}
% For the operator $A$ of Example~\ref{thm:C},
% $0\in\operatorname{conv}(\dom A)$ and
% $\operatorname{dist}(0,\dom A)\ge1/2$.
% In particular, $\overline{\dom A}$ is not convex.
% \end{corollary}
% \begin{proof}
% At $t=2$ and $t=-2$, the tails in Eq. (\ref{eq:D9}) vanish, so that
% $a^{-2}=-a^2$. Since $L$ is linear, both $La^2$ and $-La^2$ belong
% to $\dom A$. Their midpoint is zero. Eq. (\ref{eq:C9}) gives
% $\|x\|_\infty>1/2$ for every $x\in\dom A$, and hence the asserted
% distance bound. A convex closure containing $La^2$ and $-La^2$ would
% contain zero, contradicting this bound.
% \end{proof}



% \begin{remark}[Boundary points and unbounded dual values]
% \label{sec:structure}
% The nonsummable limiting block sums are also visible along a sequence
% of points in $\gra A$ whose primal components converge in norm.
% Set $t_n^\pm=\pm(1+1/n)$, $a_n^\pm=a^{t_n^\pm}$, and
% $x_n^\pm=La_n^\pm$. For $t\in J$ and $\epsilon=\sgn t$,
% direct substitution in Eqs. (\ref{eq:D4})--(\ref{eq:D9}) gives
% \[
% (La^t)_{0,1}=-2t-2\epsilon,\qquad
% (La^t)_{0,2}=-3t-2\epsilon,\qquad
% (La^t)_{0,3}=-t-\epsilon,
% \]
% with the remaining block-zero coordinates zero. In block $b\ge1$,
% the only nonzero coordinate is $-\omega_{|t|-1}(b)$ at index $1$.
% Consequently, $x_n^\pm$ converge in the supremum norm to
% $x_\pm\in c_0(I)$ with block-zero coordinates $\mp(4,5,2)$ and
% positive-block first coordinates $-1/(4b)$. Indeed, convergence of
% the tails in $\ell^2$ implies convergence in the supremum norm.
% The limits satisfy $\pair{x_\pm}{g}=\pm1$ and lie outside $\dom A$.
% On the other hand,
% \[
% \|a_n^\pm\|_1=2(1+1/n)+1+\|\omega_{1/n}\|_1\longrightarrow\infty.
% \]
% For each fixed $N$, the sum of the first $N$ tail coordinates tends to
% $\sum_{b=1}^N1/(4b)$; these partial sums are unbounded. This proves
% the divergence without asserting convergence of the labels in the dual.

% \end{remark}

\subsection{A counterexample on standard $\ell^1$}

For standard $\ell^1$, let $Z=\ell^1(\N)$ and $Z^*=\ell^\infty(\N)$, with their usual pairing. Fix once and for all an enumeration $(q_n)$ of all finitely supported rational vectors in the closed unit ball of $Y$.
\begin{definition}[The operators on standard $\ell^1$]
Define
\begin{equation}
\begin{aligned}
Qu&=\sum_{n\ge1}u_nq_n,\\
(Q^*a)_n&=\pair{q_n}{a},\\
\gra T&=\{(u,Q^*a):a\in D,\ Qu=La\},\\
f&=Q^*g,\\
Bu&=\pair{u}{f}f.
\end{aligned}
\label{eq:D10}
\end{equation}
\end{definition}
The next lemma provides the bounded preimages needed in Lemma~\ref{lem:surjection-transport}; no bounded linear right inverse of $Q$ is required.

\begin{lemma}[The specified quotient onto $c_0$]
\label{lem:specified-quotient}
The map $Q:Z=\ell^1(\N)\to Y=c_0(I)$ of Eq. (\ref{eq:D10}) is a bounded surjection with $\|Q\|=1$. Its adjoint satisfies $\|Q^*a\|_\infty=\|a\|_1$ for every $a\in\ell^1(I)$. For the fixed enumeration $(q_n)$, a map $\mathcal R:Y\to Z$ can be specified with
\[
Q\mathcal R(x)=x,\qquad \|\mathcal R(x)\|_1\le2\|x\|_\infty.
\]
\end{lemma}
\begin{proof}
\label{sec:L1}
Finitely supported rational vectors in the closed unit ball of $Y$ are countable and dense in that ball: truncate a $c_0$ vector and approximate each retained coordinate by a rational inside $[-1,1]$. Fix an enumeration of all these vectors as $(q_n)$. For $u\in Z=\ell^1(\N)$, the sum defining $Qu$ converges absolutely in the Banach space $Y$, and $\|Qu\|_\infty\le\|u\|_1$. Coordinate unit vectors occur among the $q_n$, so $\|Q\|=1$.

For an arbitrary $x\in Y$, set $r_0=x$. If $r_k\ne0$, choose the least index $n_k$ for which $q_{n_k}$ satisfies $\|q_{n_k}-r_k/\|r_k\|_\infty\|_\infty<1/2$ and set
\begin{equation}
c_k=\|r_k\|_\infty,\qquad r_{k+1}=r_k-c_kq_{n_k}.
\label{eq:L1}
\end{equation}
If a residual vanishes, stop. Otherwise $\|r_k\|_\infty\le2^{-k}\|x\|_\infty$ and $\sum_k c_k\le2\|x\|_\infty$. Define $u_n=\sum_{k:n_k=n}c_k$. Repeated indices are combined. Positivity of the $c_k$ and summability give $u\in\ell^1$ with $\|u\|_1=\sum_k c_k\le2\|x\|_\infty$. Telescoping Eq. (\ref{eq:L1}), then absolute convergence under regrouping, proves
\begin{equation}
Qu=x,\qquad \|u\|_1\le2\|x\|_\infty.
\label{eq:L2}
\end{equation}
For $x=0$, use $u=0$. This proves surjectivity with the asserted lifting bound.

\label{sec:L2}
Absolute convergence gives $\pair{u}{Q^*a}=\pair{Qu}{a}$ and $\|Q^*a\|_\infty\le\|a\|_1$. For each finite subset of $I$, its sign vector for $a$ is a finitely supported rational vector of norm at most one and occurs in the enumeration. Taking these finite sign tests proves
\begin{equation}
\|Q^*a\|_\infty=\|a\|_1\qquad(a\in\ell^1(I)).
\label{eq:L3}
\end{equation}
In particular, $Q^*$ is injective. The least-index choices in Eq. (\ref{eq:L1}) specify
\[
\mathcal R(x)=\sum_k c_k e_{n_k},\qquad \mathcal R(0)=0.
\]
The series converges in $\ell^1$ by the bound on $\sum_k c_k$, and Eq. (\ref{eq:L2}) gives both asserted properties.
\end{proof}

For this fixed enumeration,
\[
\gra T=\{(\mathcal R(La)+v,Q^*a):a\in D,\ v\in\ker Q\}.
\]
Indeed, $Q\mathcal R(La)=La$, and any other preimage differs by an element of $\ker Q$. For the fixed enumeration $(q_n)$, this formula specifies the entire graph of $T$ through the prescribed map $\mathcal R$ and the parameters $a\in D$ and $v\in\ker Q$.

% The second counterexample follows its quotient construction.
\label{sec:l1}
\begin{example}[A counterexample on standard $\ell^1$]\label{thm:L}
On $Z=\ell^1(\N)$ with its usual norm and full dual $Z^*=\ell^\infty(\N)$, let $T:Z\rightrightarrows Z^*$ and $B:Z\to Z^*$ be defined by Eq. (\ref{eq:D10}). Then:
\begin{enumerate}
\item $T$ is maximally monotone in the full pair $Z\times Z^*$.
\item $B$ is nonzero, bounded, positive, rank-one, and maximally monotone,
with $\dom B=Z$. In particular, $\dom T\cap\operatorname{int}(\dom B)\ne\varnothing$.
\item $T(0)=\varnothing$, and every $(u,b)\in\gra T$ satisfies
$\|u\|_1>1/2$ and $\pair{u}{b}\ge-2\sigma\|u\|_1$. Thus $\mathcal V(\gra T)\le2\sigma<1/4$.
\item Every $(u,c)\in\gra(T+B)$ satisfies
$\pair{u}{c}>1-\sigma>7/8$.
\item $(0,0)$ belongs to the full monotone polar of $\gra(T+B)$ and does
not belong to $\gra(T+B)$. Hence $T+B$ is not maximally monotone.
\end{enumerate}
\end{example}
\begin{proof}
Lemma~\ref{lem:specified-quotient} shows that $Q:Z\to Y$ is a bounded linear surjection for which every $x\in Y$ has a preimage of norm at most $2\|x\|_\infty$. Since $A$ is maximally monotone by Example~\ref{thm:C}, Lemma~\ref{lem:surjection-transport} applies with $C_Q=2$ and proves that $T=Q^*AQ$ is maximally monotone in $Z\times Z^*=\ell^1\times\ell^\infty$.

\noindent\textit{The radial bound on every graph value.}
\label{sec:L3}
For every $(u,Q^*a)\in\gra T$, with $t=r(a)$,
\begin{equation}
\begin{aligned}
\pair{u}{Q^*a}&=\pair{La}{a}
 =t^2-1-W(|t|-1)>-\sigma,\\
\|u\|_1&\ge\|Qu\|_\infty=\|La\|_\infty>1/2.
\end{aligned}
\label{eq:L6}
\end{equation}
Thus $T(0)=\varnothing$ and $\pair{u}{Q^*a}\ge-2\sigma\|u\|_1$ on the entire graph, including every $\ker Q$ translate. Consequently, $\mathcal V(\gra T)\le2\sigma<1/4$.

\noindent\textit{The operator $B$ and the domain intersection condition.}
\label{sec:L4}
Let $f=Q^*g$. By Eq. (\ref{eq:L3}), $\|f\|_\infty=\|g\|_1=2$, so $f$ is nonzero. The map $B:u\mapsto\pair{u}{f}f$ is nonzero, linear, bounded with norm at most $4$, positive, and rank-one. The rank-one proof in Theorem~\ref{thm:endpoint-criterion}, now testing all $v\in\ell^1$ and dual elements in $\ell^\infty$, proves $B$ maximally monotone. Its domain is all of $Z$, so $\operatorname{int}(\dom B)=Z$.

The vector $q_\circ=-(3/4)e_{0,1}-e_{0,2}-(3/8)e_{0,3}$ occurs as $q_{n_0}$. For $u_\circ=8e_{n_0}$, $Qu_\circ=La^2$. Thus $(u_\circ,Q^*a^2)\in\gra T$, and $\dom T\cap\operatorname{int}(\dom B)\ne\varnothing$.

\noindent\textit{Nonmaximality of $T+B$.}
\label{sec:L5}
For every $(u,Q^*a)\in\gra T$, let $t=r(a)$. Then $\pair{u}{f}=\pair{Qu}{g}=\pair{La}{g}=t$. Consequently, for every $(u,Q^*a)\in\gra T$,
\begin{equation}
\pair{u}{Q^*a+Bu}
 =2t^2-1-W(|t|-1)>1-\sigma>7/8.
\label{eq:L7}
\end{equation}
The sum is monotone by monotonicity of its factors. Eq. (\ref{eq:L7}) shows that $(0,0)$ is related to every point of its graph. Since $T(0)=\varnothing$, $(0,0)\notin\gra(T+B)$. Adjoining this point gives a proper monotone extension of $\gra(T+B)$.
\end{proof}

\section{Conclusions}
\label{sec:conclusion}

We gave counterexamples to Rockafellar's sum conjecture on $c_0$ and standard $\ell^1$. Our structural theorem computes the entire monotone polar for the stated class of graphs and identifies exactly which equations in $X^*$ determine maximality. In the coordinate construction, solutions with parameter $\tau\in[-1,1]$ would require a nonsummable sequence of block sums. The pairing identity proves monotonicity and the finite radial bound, and shows why adding a full-domain positive rank-one operator makes the sum nonmaximal. The pullback through a bounded surjection gives the second counterexample, with maximality verified in the full $\ell^1\times\ell^\infty$ pair.


\section*{Use of AI}
The author supplied earlier constructions and results, involving the sign function and triangular operators, together with a framework for positive rank-one perturbations, and guided their further development to investigate how maximality can fail under addition. The author also supplied earlier unsuccessful constructions and obstruction results showing that a finite radial bound on a monotone graph need not survive maximal extension. Through iterative discussions of these materials, the author and GPT-5.6 Sol developed a general construction theorem characterizing maximality and identifying when an everywhere-defined rank-one monotone perturbation yields a nonmaximal sum.  Under the direction of the author and the aforementioned research results, GPT-5.6 Sol worked out the detailed constructions and proof arguments, including the coupling of triangular operators with a Lipschitz curve whose limits at $t=-1$ and $t=1$ lie outside the range of the block-sum map on $(c_0)^*$, and the transfer from $c_0$ to standard $\ell^1$. 





% and directed the counterexample search and a general construction theorem characterizing maximality and showing when addition of an everywhere-defined rank-one monotone operator produces a nonmaximal sum, and. Under this direction, GPT-5.6 Sol developed explicit formulas and proof arguments for the counterexamples on $c_0$ and standard $\ell^1$, and assisted with literature checks, adversarial review and exposition. 

\begingroup
\raggedright \bibliographystyle{IEEEtran}
\bibliography{refer}
\endgroup

\end{document}
